\documentclass[a4paper, 11pt]{article}

\usepackage[margin=3cm]{geometry} 
\usepackage{amsthm, amsmath, amssymb}
\usepackage{array}
\usepackage{bussproofs}
\usepackage{booktabs}

\theoremstyle{definition}
\newtheorem{definition}{Definition}[section]
\newtheorem{example}[definition]{Example}
\newtheorem{theorem}[definition]{Theorem}
\newtheorem{lemma}[definition]{Lemma}
\newtheorem{prop}[definition]{Proposition}
\newtheorem{corollary}[definition]{Corollary}

% change spacing
\let\olditemize\itemize
\renewcommand{\itemize}{\olditemize\setlength{\itemsep}{0pt}\setlength{\parsep}{0pt}}
\let\oldenumerate\enumerate
\renewcommand{\enumerate}{\oldenumerate\setlength{\itemsep}{0pt}\setlength{\parsep}{0pt}}

% hilbert proofs
\newcounter{dline}
\newenvironment{deduction}[1][\textit{no hypotheses}]
    {\setcounter{dline}{0}%
     \[\begin{array}{r|l@{\qquad}l}
     & #1 \\ \cline{2-2}}
    {\end{array}\]}
\newcommand{\dl}[2]{\stepcounter{dline}\thedline & #1 & \text{#2} \\}

% tables
\newcounter{factrow}
\newenvironment{comparison}[4][p{0.22\textwidth} p{0.36\textwidth} p{0.36\textwidth}]
    {\setcounter{factrow}{0}%
     \begin{center}
     \begin{tabular}{#1}
     \textsc{#2} & \textsc{#3} & \textsc{#4} \\[4pt]}
    {\end{tabular}
     \end{center}}

% ======================================================================
% WORKSHEET MACROS: you should only need the commands in this cheat sheet.
%
%   SCRIPTS (spaces are application, so just type them as on the sheet)
%     \s{S K K a}             inline script, e.g. \s{(\lam{x} x z) y}
%     \lam{x}                 lambda x.  (use inside scripts)
%
%   RUNS (each \st is one step; the step count is printed for you)
%     \begin{run}
%       \from{K a b c}        the starting script (already filled in)
%       \st{a c}              one step
%       \st{a c}[Int]         ... optionally with a type
%     \end{run}
%     Write \st{...} with no space before the brace.
%
%   CERTIFICATES (numbered automatically)
%     \begin{deduction}[p \to q,\ p]       hypotheses (leave off for none)
%       \dl{p \to q}{hypothesis}
%       \dl{q \to p \to q}{\axI{q, p}}     (I) [q, p]
%       \dl{q}{\MP{1}{2}}                  MP 1, 2
%       \ddl{f x}{q}{\MP{1}{2}}            decorated line: script : statement
%     \end{deduction}
%
%   TYPINGS
%     \begin{typing}
%       \ty{K K}{C \to A \to B \to A}{(app) 1, 2}   prints  n. |- K K : ...
%     \end{typing}
%
%   TREES (write the tree bottom-up in reading order: premises first, then
%   the bar that joins them; bussproofs does the layout)
%     \begin{prooftree}
%       \hyp{A \to B}  \dis{A}{1}     open leaf, then leaf [A]^1
%       \impE{B}                      bar with 2 premises
%       \impI[1]{A \to B}             bar discharging 1  (\impI{..} = none)
%     \end{prooftree}
%     Bars: \impE \impI[n] \andI \andEi \andEii \orIi \orIii \orE{i}{j}
%           \botE \allE \allI \exI \exE{n} \eqI \eqE
%     Anything else: \unary{label}{A} \binary{label}{A} \trinary{label}{A}
%     \goal{A}   placeholder: shows A under a \hyp{\vdots}; delete when done
%
%   GAPS   \fillin{}  prints a blank line; type inside the braces to fill it
% ======================================================================
\usepackage{xcolor}
\definecolor{qgrey}{gray}{0.4}
\makeatletter

% questions: the sheet's text in grey, your answer below it
\newenvironment{q}[1]
    {\subsection*{#1}\begingroup\small\color{qgrey}}
    {\par\endgroup\medskip}
\newenvironment{parts}
    {\begin{enumerate}\renewcommand{\labelenumi}{(\roman{enumi})}}
    {\end{enumerate}}
\newcommand{\fillin}[1]{\if\relax\detokenize{#1}\relax\rule[-2pt]{3em}{0.4pt}\else#1\fi}

% scripts: inside \s{..}, \from, \st, \ty and \ddl a typed space is application
\begingroup\catcode`\ =\active%
\gdef {\ifmmode\ \else\space\fi}%
\endgroup
\newcommand{\appspaces}{\catcode`\ =\active}
\newcommand{\s}{\begingroup\appspaces\s@}
\newcommand{\s@}[1]{\endgroup\ensuremath{#1}}
\newcommand{\lam}[1]{\lambda #1.}

% runs
\newcounter{rstep}
\newenvironment{run}
    {\setcounter{rstep}{0}\[\begin{array}{r@{\quad}l@{\qquad}l}}
    {& \ifnum\value{rstep}>0 \text{\small\itshape (\arabic{rstep}~step\ifnum\value{rstep}=1\else s\fi)}\fi \end{array}\]}
\newcommand{\from}{\begingroup\appspaces\from@}
\newcommand{\from@}[1]{\endgroup\@ifnextchar[{\run@line{}{#1}}{\run@line{}{#1}[]}}
\newcommand{\st}{\begingroup\appspaces\st@}
\newcommand{\st@}[1]{\endgroup\stepcounter{rstep}%
    \@ifnextchar[{\run@line{\rightsquigarrow}{#1}}{\run@line{\rightsquigarrow}{#1}[]}}
\def\run@line#1#2[#3]{#1 & #2 & \if\relax\detokenize{#3}\relax\else : #3\fi \\}

% certificates and typings
\newcommand{\axI}[1]{(I)~$[#1]$}
\newcommand{\axII}[1]{(II)~$[#1]$}
\newcommand{\MP}[2]{MP~#1,~#2}
\newcommand{\ddl}{\begingroup\appspaces\ddl@}
\newcommand{\ddl@}[1]{\endgroup\ddl@@{#1}}
\newcommand{\ddl@@}[3]{\dl{#1 : #2}{#3}}
\newenvironment{typing}
    {\setcounter{dline}{0}\[\begin{array}{r@{.\ \ }l@{\qquad}l}}
    {\end{array}\]}
\newcommand{\ty}{\begingroup\appspaces\ty@}
\newcommand{\ty@}[1]{\endgroup\ty@@{#1}}
\newcommand{\ty@@}[3]{\stepcounter{dline}\thedline & \vdash #1 : #2 & \text{#3} \\}

% trees
\newcommand{\hyp}[1]{\AxiomC{$#1$}}
\newcommand{\dis}[2]{\AxiomC{$[#1]^{#2}$}}
\newcommand{\goal}[1]{\noLine\UnaryInfC{$#1$}}
\newcommand{\rl}[2]{\RightLabel{\scriptsize #1\if\relax\detokenize{#2}\relax\else, #2\fi}}
\newcommand{\unary}[2]{\RightLabel{\scriptsize #1}\UnaryInfC{$#2$}}
\newcommand{\binary}[2]{\RightLabel{\scriptsize #1}\BinaryInfC{$#2$}}
\newcommand{\trinary}[2]{\RightLabel{\scriptsize #1}\TrinaryInfC{$#2$}}
\newcommand{\impE}[1]{\binary{$\to$E}{#1}}
\newcommand{\impI}[2][]{\rl{$\to$I}{#1}\UnaryInfC{$#2$}}
\newcommand{\andI}[1]{\binary{$\land$I}{#1}}
\newcommand{\andEi}[1]{\unary{$\land$E$_1$}{#1}}
\newcommand{\andEii}[1]{\unary{$\land$E$_2$}{#1}}
\newcommand{\orIi}[1]{\unary{$\lor$I$_1$}{#1}}
\newcommand{\orIii}[1]{\unary{$\lor$I$_2$}{#1}}
\newcommand{\orE}[3]{\trinary{$\lor$E, #1, #2}{#3}}
\newcommand{\botE}[1]{\unary{$\bot$E}{#1}}
\newcommand{\allE}[1]{\unary{$\forall$E}{#1}}
\newcommand{\allI}[1]{\unary{$\forall$I}{#1}}
\newcommand{\exI}[1]{\unary{$\exists$I}{#1}}
\newcommand{\exE}[2]{\binary{$\exists$E, #1}{#2}}
\newcommand{\eqI}[1]{\AxiomC{}\unary{$=$I}{#1}}
\newcommand{\eqE}[1]{\binary{$=$E}{#1}}

% names
\newcommand{\nm}[1]{\mathit{#1}}
\newcommand{\str}[1]{\text{\texttt{"#1"}}}
\newcommand{\Int}{\mathsf{Int}}
\newcommand{\String}{\mathsf{String}}
\newcommand{\fst}{\mathsf{fst}}
\newcommand{\snd}{\mathsf{snd}}
\newcommand{\Left}{\mathsf{Left}}
\newcommand{\Right}{\mathsf{Right}}
\newcommand{\absurd}{\mathsf{absurd}}
\newcommand{\caseof}[5]{\mathsf{case}\ #1\ \mathsf{of}\ \Left\ #2 \Rightarrow #3 \mid \Right\ #4 \Rightarrow #5}
\makeatother

\begin{document}

\section*{Problem Sheet 2 --- Part B}

% ---------------------------------------------------------------------
\begin{q}{B1. Brackets}
Insert the brackets that right-grouping hides: (i) $p \to q \to r \to s$; (ii) $(p \to q) \to r \to s$;
(iii) $p \to (q \to r) \to s$. Then say which two of these are the same statement: (iv) $p \to (q \to r)$;
(v) $(p \to q) \to r$; (vi) $p \to q \to r$.
\end{q}
\begin{parts}
    \item $p \to q \to r \to s$ \quad = \quad $ $
    \item $(p \to q) \to r \to s$ \quad = \quad $ $
    \item $p \to (q \to r) \to s$ \quad = \quad $ $
\end{parts}
The same statement: 

% ---------------------------------------------------------------------
\begin{q}{B2. Name the axiom}
Fill in the reason for each line: (I) or (II) with what was put in place of the letters, or ``neither''.
\end{q}
\begin{parts}
    \item $q \to p \to q$: 
    \item $(p \to q) \to (p \to q) \to (p \to q)$: 
    \item $p \to (q \to r) \to p$: 
    \item $(p \to q) \to p \to q$: 
    \item $(q \to q \to q) \to (q \to q) \to q \to q$: 
    \item $(p \to q \to r) \to (p \to r) \to p \to r$: 
\end{parts}

% ---------------------------------------------------------------------
\begin{q}{B3. Fill in the statement}
Fill in the statement for each reason. Then fill in the gaps in the certificate from the hypothesis $p \to q \to r$.
\end{q}
\begin{parts}
    \item \axI{p, p}: \quad $ $
    \item \axI{q \to p, r}: \quad $ $
    \item \axII{p, p \to q, r}: \quad $ $
\end{parts}
\begin{deduction}[p \to q \to r]
    \dl{p \to q \to r}{hypothesis}
    \dl{\fillin{}}{\axII{p, q, r}}
    \dl{\fillin{}}{\MP{2}{1}}
    \dl{\fillin{}}{\axI{(p \to q) \to p \to r,\ q}}
    \dl{\fillin{}}{\MP{4}{3}}
\end{deduction}

% ---------------------------------------------------------------------
\begin{q}{B4. Is it a certificate?}
Is each list a certificate for its last line from the hypotheses given? If not, find the first line that fails, and say why.
\end{q}
\begin{parts}
    \item \begin{deduction}[p \to q,\ q \to r,\ p]
        \dl{p}{hypothesis}
        \dl{p \to q}{hypothesis}
        \dl{q}{\MP{2}{1}}
        \dl{q \to r}{hypothesis}
        \dl{r}{\MP{4}{3}}
    \end{deduction}
    Answer: 
    \item \begin{deduction}[p \to q,\ q]
        \dl{p \to q}{hypothesis}
        \dl{q}{hypothesis}
        \dl{p}{\MP{1}{2}}
    \end{deduction}
    Answer: 
    \item \begin{deduction}[q]
        \dl{q}{hypothesis}
        \dl{q \to p \to q}{\axI{q, p}}
        \dl{p \to q}{\MP{2}{1}}
    \end{deduction}
    Answer: 
    \item \begin{deduction}[p \to q,\ p]
        \dl{p \to q}{hypothesis}
        \dl{q}{\MP{1}{3}}
        \dl{p}{hypothesis}
    \end{deduction}
    Answer: 
\end{parts}

% ---------------------------------------------------------------------
\begin{q}{B5. Write a certificate}
(i) $p \to q \to r,\ p,\ q \vdash r$\quad (ii) $p \to q,\ q \to r,\ p \vdash r$\quad (iii) $q \vdash p \to q$\quad
(iv) $p \to q \to r,\ q \vdash p \to r$ (seven lines; you need one (I) and one (II) line)
\end{q}
\begin{parts}
    \item \begin{deduction}[p \to q \to r,\ p,\ q]
        \dl{ }{ }
    \end{deduction}
    \item \begin{deduction}[p \to q,\ q \to r,\ p]
        \dl{ }{ }
    \end{deduction}
    \item \begin{deduction}[q]
        \dl{ }{ }
    \end{deduction}
    \item \begin{deduction}[p \to q \to r,\ q]
        \dl{ }{ }
    \end{deduction}
\end{parts}

% ---------------------------------------------------------------------
\begin{q}{B6. Gluing certificates}
A2(b)'s certificate, with $p, q, r$ for $y, s, m$, shows $p \to q,\ q \to r \vdash p \to r$ in seven lines; with $r, s$ for
$q, r$ it shows $p \to r,\ r \to s \vdash p \to s$.
(i) Explain how to glue the two into a certificate for $p \to q,\ q \to r,\ r \to s \vdash p \to s$, and say how many lines
it has. (ii) Why is a certificate from $\Gamma$ also a certificate from $\Gamma$ with one more hypothesis added?
(iii) Can a certificate use infinitely many hypotheses? Why not?
\end{q}
\begin{parts}
    \item 
    \item 
    \item 
\end{parts}

% ---------------------------------------------------------------------
\begin{q}{B7. Rub out the scripts}
On Sheet 1, B11 typed \s{K K} like this:
\begin{typing}
    \ty{K}{(A \to B \to A) \to C \to A \to B \to A}{(K)}
    \ty{K}{A \to B \to A}{(K)}
    \ty{K K}{C \to A \to B \to A}{(app) 1, 2}
\end{typing}
Erase it, with $p, q, r$ for $A, B, C$, and give the justification of each line. Then run \s{K K x y z} and say what the
script does.
\end{q}
\begin{deduction}
    \dl{ }{ }
    \dl{ }{ }
    \dl{ }{ }
\end{deduction}
\begin{run}
    \from{K K x y z}
\end{run}
It does: 

% ---------------------------------------------------------------------
\begin{q}{B8. Flip, without the recipe}
Here is a certificate for $p \to q \to r \vdash q \to p \to r$. Check lines 4 and 6. Decorate the certificate, calling
the hypothesis $f$, and run the script on line 9 applied to $b$ and then $a$. What does it do?
\end{q}
\begin{deduction}[p \to q \to r]
    \dl{p \to q \to r}{hypothesis}
    \dl{(p \to q \to r) \to (p \to q) \to p \to r}{\axII{p, q, r}}
    \dl{(p \to q) \to p \to r}{\MP{2}{1}}
    \dl{((p \to q) \to p \to r) \to q \to (p \to q) \to p \to r}{\axI{(p \to q) \to p \to r,\ q}}
    \dl{q \to (p \to q) \to p \to r}{\MP{4}{3}}
    \dl{(q \to (p \to q) \to p \to r) \to (q \to p \to q) \to q \to p \to r}{\axII{q,\ p \to q,\ p \to r}}
    \dl{(q \to p \to q) \to q \to p \to r}{\MP{6}{5}}
    \dl{q \to p \to q}{\axI{q, p}}
    \dl{q \to p \to r}{\MP{7}{8}}
\end{deduction}
Lines 4 and 6: 

Decorated (fill in each script):
\begin{deduction}[f : p \to q \to r]
    \ddl{f}{p \to q \to r}{hypothesis}
    \ddl{S}{(p \to q \to r) \to (p \to q) \to p \to r}{\axII{p, q, r}}
    \ddl{ }{(p \to q) \to p \to r}{\MP{2}{1}}
    \ddl{ }{((p \to q) \to p \to r) \to q \to (p \to q) \to p \to r}{\axI{(p \to q) \to p \to r,\ q}}
    \ddl{ }{q \to (p \to q) \to p \to r}{\MP{4}{3}}
    \ddl{ }{(q \to (p \to q) \to p \to r) \to (q \to p \to q) \to q \to p \to r}{\axII{q,\ p \to q,\ p \to r}}
    \ddl{ }{(q \to p \to q) \to q \to p \to r}{\MP{6}{5}}
    \ddl{ }{q \to p \to q}{\axI{q, p}}
    \ddl{ }{q \to p \to r}{\MP{7}{8}}
\end{deduction}
\begin{run}
    \from{ b a}   % put line 9's script in front of b a
\end{run}
It does: 

% ---------------------------------------------------------------------
\begin{q}{B9. Two detours}
Decorate this certificate for $q$ from $p \to q,\ p,\ r,\ s$, calling the hypotheses $f, x, y, z$, and run the script on
its last line to the end. How many lines does the certificate for the final script have?
\end{q}
\begin{deduction}[f : p \to q,\ x : p,\ y : r,\ z : s]
    \ddl{f}{p \to q}{hypothesis}
    \ddl{x}{p}{hypothesis}
    \ddl{ }{q}{\MP{1}{2}}
    \ddl{ }{q \to r \to q}{\axI{q, r}}
    \ddl{ }{r \to q}{\MP{4}{3}}
    \ddl{y}{r}{hypothesis}
    \ddl{ }{q}{\MP{5}{6}}
    \ddl{ }{q \to s \to q}{\axI{q, s}}
    \ddl{ }{s \to q}{\MP{8}{7}}
    \ddl{z}{s}{hypothesis}
    \ddl{ }{q}{\MP{9}{10}}
\end{deduction}
\begin{run}
    \from{ }   % line 11's script
\end{run}
Lines in the certificate for the final script: 

% ---------------------------------------------------------------------
\begin{q}{B10. Use the recipe of A3}
(i) $\vdash p \to p$. Which lines does the machine write for you? (ii) $\vdash (p \to q) \to (q \to r) \to p \to r$
(iii) $\vdash p \to (p \to q) \to q$ (iv) $p \to q \to r \vdash q \to p \to r$ (compare B8)
\end{q}
\begin{parts}
    \item \begin{deduction}
        \dl{ }{ }
    \end{deduction}
    The machine writes: 
    \item \begin{deduction}
        \dl{ }{ }
    \end{deduction}
    \item \begin{deduction}
        \dl{ }{ }
    \end{deduction}
    \item \begin{deduction}[p \to q \to r]
        \dl{ }{ }
    \end{deduction}
\end{parts}

% ---------------------------------------------------------------------
\begin{q}{B11. Run the machine}
Run the machine of A3(c) on this certificate for $p \vdash q \to p$, discharging $p$. How many lines does the machine
write, and what is the last line? Is there a shorter certificate for that? Decorate the machine's output and run the
script on $x$ and then $y$.
\begin{deduction}[p]
    \dl{p \to q \to p}{\axI{p, q}}
    \dl{p}{hypothesis}
    \dl{q \to p}{\MP{1}{2}}
\end{deduction}
\end{q}
\begin{deduction}
    \ddl{ }{ }{ }
\end{deduction}
Number of lines: \qquad Last line: 

Shorter certificate: 
\begin{run}
    \from{ x y}   % put the last line's script in front of x y
\end{run}

% ---------------------------------------------------------------------
\begin{q}{B12. Counting letters}
In A3, line 1 of Alice's certificate was (II) $[P, Q, R]$ with $P$ for $s \to m$, $Q$ for $y \to s \to m$ and $R$ for
$(y \to s) \to y \to m$. Write it out with $P, Q, R$ replaced, and confirm that it has twenty letters and nineteen arrows.
Line 1 redoes, one level up and under the hypothesis $P$, one of the MP lines of A2(b)'s certificate. Which one?
\end{q}
\[ \]
Letters: \qquad Arrows: 

Which MP line: 

% ---------------------------------------------------------------------
\begin{q}{B13. A shortcut for the machine (for the curious)}
Suppose a line $C$ was obtained without using $A$ anywhere above it. Show that, if the machine keeps the lines that obtained
$C$, two more lines give $A \to C$, whatever kind of line $C$ is. What does this shortcut do to scripts? Why does it make
the machine's output shorter?
\end{q}
\begin{deduction}[\ldots]
    \dl{C}{(however $C$ was obtained)}
    \dl{ }{ }
    \dl{ }{ }
\end{deduction}
To scripts: 

Why shorter: 

% ---------------------------------------------------------------------
\begin{q}{B14. Write a checker (for the curious)}
Write a checker: a program (in Python, Haskell, or careful pseudocode) that takes a list of statements, each with its
justification, and a list of hypotheses, and says whether the list is a certificate. How long does it take, in the length
of the list? Which part of your program corresponds to rule (var) of a type checker, and which to (app)?
\end{q}
\begin{verbatim}
def is_certificate(lines, hypotheses):
    ...
\end{verbatim}
Running time: 

(var) / (app): 

\end{document}
